% Generated draft: requires mathematical and source review.
\documentclass[11pt]{article}
\usepackage[margin=25mm]{geometry}
\usepackage{fontspec}
\setmainfont{Latin Modern Roman}
\usepackage{amsmath,amssymb,mathtools,unicode-math}
\setmathfont{Latin Modern Math}
\usepackage{enumitem,needspace,xurl,hyperref}
\setlength{\parindent}{0pt}
\setlength{\parskip}{5pt}
\setlength{\emergencystretch}{3em}
\newcommand{\N}{\mathbb{N}}
\newcommand{\Z}{\mathbb{Z}}
\newcommand{\Q}{\mathbb{Q}}
\newcommand{\R}{\mathbb{R}}
\newcommand{\C}{\mathbb{C}}
\newcommand{\F}{\mathbb{F}}
\newcommand{\A}{\mathbb{A}}
\newcommand{\PP}{\mathbb{P}}
\newcommand{\E}{\mathbb{E}}
\newcommand{\eps}{\varepsilon}
\newcommand{\dd}{\,\mathrm{d}}
\begin{document}
\section*{UnsolvedMath ID: 701506}
\textbf{Source problem code:} catalogue entry
\subsection{Definitions and mathematical statement}
Let \(p\) be a prime and let \(G\) be a finite group of order \(p^n\) for some integer \(n\geq 0\). An irreducible complex representation of \(G\) is a homomorphism \(\Theta:G\to\operatorname{GL}(V_\Theta)\), where \(V_\Theta\) is a nonzero finite-dimensional complex vector space with no proper nonzero \(G\)-invariant subspace. Its degree is \(\dim_{\C}V_\Theta\), and its kernel is \(\ker\Theta=\{g\in G:\Theta(g)=\operatorname{id}_{V_\Theta}\}\).

Let \(e\geq 0\) be the integer such that \(p^e\) is the largest degree of an irreducible complex representation of \(G\), and let \(\mathcal M(G)\) contain one representative from each equivalence class of irreducible complex representations of degree \(p^e\). Is it true, for every such prime \(p\) and finite \(p\)-group \(G\), that
\[
 p>e\quad\Longrightarrow\quad
 \bigcap_{\Theta\in\mathcal M(G)}\ker\Theta=\{1_G\}?
\]
The supplied registry labels this question open; that status is not independently verified here.
\subsection{Short English statement}
If the largest irreducible complex representation degree of a finite \(p\)-group is \(p^e\) and \(p>e\), must every nonidentity element act nontrivially in at least one irreducible representation of that largest degree?
\subsection{Sources}
\begin{itemize}
\item Local UnsolvedMath registry, supplied source record: permanent ID \(701506\), exact source code \texttt{catalogue entry}, title \emph{Kourovka Notebook Problem 18.77}. The supplied status is \texttt{open}.
\item \emph{The Kourovka Notebook (Unsolved Problems in Group Theory)}, Problem 18.77, proposed by D. S. Passman in 2014. The supplied background identifies the 18th issue (2014) and the 21st issue edition (2026), arXiv:1401.0300v48. The supplied collection description names E. I. Khukhro and V. D. Mazurov as editors. This bibliographic information is supplied attribution and has not been independently verified.
\end{itemize}
\end{document}
